\subsection{The weight tail in Beck and special PD linearization}
\label{arch136:tail-conormal}

\paragraph{ARCH-136-TAIL-CONORMAL: scope.}
This is an explicit calculation for this investigation, with no historical
priority assertion.
This calculation distinguishes the conormal module for PD Beck modules
from the conormal module for special PD derivations. It does not compute
the full derived cotangent complex or assert that a simultaneous
admissible completion exists. In particular, the infinite projective
dimension in the first calculation is not an obstruction in the special
differential category used in cote~136.

Let \(R=\mathbb Z_{(p)}\), \(S=\Gamma_R(T)\), and
\(I=I_{\ge N}=\bigoplus_{w\ge N}R\,T^{[w]}\).
Write \(A=S/I\), \(\mathfrak a=\ker(A\to R)\), and
\(M=I/I^2\). The square-zero PD ideal \(M\) has its canonical structure
as a PD Beck module over the nonunital PD algebra \(\mathfrak a\).
Over \(R\), the operator algebra of an abelian PD algebra is
\[
 B=U_R(0)=R[\phi]/(p\phi),
 \qquad
 U=U_R(\mathfrak a)=A\otimes_R B
 \tag{CON.1}
\]
as an \(R\)-module. Its multiplication is determined by
\(\phi a=0\) for \(a\in\mathfrak a\), together with the algebra
multiplications on \(A\) and \(B\). Here \(\phi\) records \(\gamma_p\).
The other prime operators vanish after localization, and the scalar
Frobenius twist is ordinary scalar multiplication since
\(r^p-r\in pR\) and \(p\phi=0\).
These are the specialization of Kmail--Kozak--Miller,
\emph{Divided powers and K\"ahler differentials},
Communications in Algebra \textbf{54}(5) (2026), 2285--2293,
Proposition~4.2, to \(R=\mathbb Z_{(p)}\).
The author's version gives this proposition on printed pp.~5--6;
the DOI is \texttt{10.1080/00927872.2025.2578689}.

Define the special conormal by
\[
 M_{\mathrm{sp}}
   =M/\bigl(A\cdot\phi(M)\bigr)
   =I/I^{[2]},
 \tag{CON.2}
\]
where \(I^{[2]}\) is the ideal generated by \(I^2\) and
\(\gamma_r(x)\), \(x\in I,\ r\ge2\). The \(A\)-span in (CON.2) is
essential; the set-theoretic image of \(\phi\) need not be an
\(A\)-submodule.

\paragraph{Large-prime calculation.}
Suppose \(N\ge2\) and \(p>2N\). Then \(A\cong R[t]/(t^N)\), and
\[
 M\cong A e\ \oplus\!
       \bigoplus_{i\ge1}(A/p)v_i,
 \quad
 \operatorname{wt}(e)=N,\quad \operatorname{wt}(v_i)=p^i.
 \tag{CON.3}
\]
The generators \(v_i\) can be chosen so that
\(\phi(e)=0\) and \(\phi(v_i)=v_{i+1}\).
Consequently
\[
 \operatorname{pd}_{U}(M)=\infty,
 \qquad
 M_{\mathrm{sp}}\cong A e\oplus(A/p)v_1,
 \qquad
 \operatorname{pd}_{A}(M_{\mathrm{sp}})=1.
 \tag{CON.4}
\]

\emph{Proof.}
Put \(x_i=T^{[p^i]}\). For units \(u_i\in R^\times\),
\[
 S\cong
 R[x_0,x_1,\ldots]/(x_i^p-pu_i x_{i+1}:i\ge0).
 \tag{CON.5}
\]
The digit monomials \(\prod x_i^{e_i}\), \(0\le e_i<p\), map to
unit multiples of the divided-power basis. Reducing by the displayed
monic relations, separately in each weight, therefore proves (CON.5).
As \(N<p\),
\(I=(x_0^N,x_1,x_2,\ldots)\), and the quotient is \(R[t]/(t^N)\).

Let \(P=R[x_0,x_1,\ldots]\), let \(F\) be the ideal of relations in
(CON.5), and put \(J=(x_0^N,x_1,x_2,\ldots)\).
Then \(F\subseteq J\), \(P/J=A\), and
\[
 M=J/(J^2+F).
\]
The conormal \(J/J^2\) is free over \(A\), with generators
\(e=x_0^N\) and \(v_i=x_i\), \(i\ge1\). Since \(p>2N\),
\(x_0^p\in J^2\); also \(x_i^p\in J^2\) for \(i\ge1\).
The images of the defining relations in \(J/J^2\) are thus precisely
\(-pu_i v_{i+1}\). This proves the underlying module formula (CON.3).

In \(M\), the operation \(\phi\) is induced by \(\gamma_p\). We have
\[
 \gamma_p(x_0^N)=x_0^{p(N-1)}\gamma_p(x_0)
                =x_0^{p(N-1)}x_1\in I^2,
\]
so \(\phi(e)=0\). Also
\[
 \gamma_p(T^{[p^i]})
   =\frac{(p^{i+1})!}{p!(p^i!)^p}\,T^{[p^{i+1}]},
\]
whose coefficient is a \(p\)-unit. Rescaling the \(v_i\) by units
modulo \(p\), recursively, makes \(\phi(v_i)=v_{i+1}\).
The operator annihilates \(\mathfrak a M\), as required in (CON.1).
This also follows directly from the product formula and \(t^p=0\)
in \(A\).

Let
\[
 V=B/(\phi)\oplus B/(p)=R\oplus\mathbb F_p[\phi].
\]
The preceding action formula identifies \(M\), including its Beck
module structure, with \(U\otimes_B V\). Here \(U\) is free as a
right \(B\)-module, since \(A\) is \(R\)-free. The \(B\)-module
\(R=B/(\phi)\) has the exact free resolution with alternating
differentials \(\phi,p\):
\[
 \cdots\xrightarrow{\phi}B\xrightarrow{p}B
 \xrightarrow{\phi}B\longrightarrow R\longrightarrow0.
 \tag{CON.6}
\]
Indeed \(\operatorname{ann}(\phi)=(p)\) and
\(\operatorname{ann}(p)=(\phi)\). Tensoring with \(R\) shows
\[
 \operatorname{Tor}_{2j+1}^B(R,R)\cong\mathbb F_p
 \quad(j\ge0).
\]
Thus \(V\), which contains \(R\) as a direct summand, has infinite
\(B\)-projective dimension.

The augmentation \(U\to B\) kills \(\mathfrak a\). To compute
\(\operatorname{Tor}^U(B,M)\), induce a projective \(B\)-resolution
of \(V\) to \(U\). This remains a resolution because \(U\) is
right-\(B\)-free, and base change back to \(B\) returns the original
resolution. Hence
\[
 B\otimes_U M\cong V,\qquad
 \operatorname{Tor}_i^U(B,M)=0\quad(i>0).
\]
If \(M\) had a finite \(U\)-projective resolution, its base change
would therefore be an exact finite \(B\)-projective resolution of
\(V\), a contradiction.

Finally the \(A\)-submodule generated by \(\phi(M)\) is exactly the
sum of the \(v_i\)-summands with \(i\ge2\).
This gives the special quotient in (CON.4), with free resolution
\[
 0\longrightarrow A
   \xrightarrow{\ r\mapsto(0,pr)\ }
   A\oplus A\longrightarrow M_{\mathrm{sp}}\longrightarrow0.
 \tag{CON.7}
\]
For the weight grading, the first copy of \(A\) in the middle is
shifted by \(N\), and the other two copies are shifted by \(p\).
Injectivity follows from the \(R\)-freeness of \(A\).
The nonzero summand \(A/p\) is not \(A\)-projective: any
\(A\)-projective is \(R\)-flat, whereas \(A/p\) is not.
Thus the projective dimension in (CON.7) is exactly one.
\(\square\)

\paragraph{Uniform special conormal after a finite localization.}
For \(N\ge2\), put \(R_0=\mathbb Z[1/(2N)!]\),
\(S_0=\Gamma_{R_0}(T)\), \(I_0=I_{\ge N}\), and
\(A_0=S_0/I_0\cong R_0[t]/(t^N)\). Then
\[
 I_0/I_0^{[2]}
 \cong A_0e_N\oplus
       \bigoplus_{\substack{p\ {\rm prime}\\p>2N}}(A_0/p)v_p,
 \quad \operatorname{wt}(e_N)=N,\quad\operatorname{wt}(v_p)=p.
 \tag{CON.8}
\]
In particular this module has the single weight-graded free resolution
\[
 0\longrightarrow
 \bigoplus_{p>2N}A_0r_p
 \xrightarrow{\ r_p\mapsto p v_p\ }
 A_0e_N\oplus\bigoplus_{p>2N}A_0v_p
 \longrightarrow I_0/I_0^{[2]}\longrightarrow0,
 \quad \operatorname{wt}(r_p)=p.
 \tag{CON.9}
\]

\emph{Proof.}
The submodule generated by \(T^{[N]}\) is free over \(A_0\):
it occupies exactly weights \(N,\ldots,2N-1\), and the relevant
binomial coefficients are units in \(R_0\).
Every remaining weight is torsion, since rationally
\(I_0/I_0^{[2]}\) is \((T^N)/(T^{2N})\).
Each weight component is a quotient of the rank-one \(R_0\)-module
\(R_0T^{[w]}\), so its torsion has the usual finite primary
decomposition. Localization of PD powers and of the quotient commutes
with these operations. For each retained prime \(p>2N\), (CON.4)
says that the \(p\)-primary summand is killed by \(p\), has weights
\(p,\ldots,p+N-1\), and is free of rank one over \(A_0/p\).
Take \(v_p\) to be the \(p\)-primary component of the class of
\(T^{[p]}\). This generates that summand by the localized formula.
It is important to take the primary component: the raw class of
\(T^{[p]}\) can also have torsion at a nearby smaller prime.
The action of \(A_0\) preserves each primary summand, proving (CON.8).
Multiplication by \(p\) is injective on \(A_0\), which proves the
exactness of (CON.9). \(\square\)

\paragraph{The \(p\)-power contrast.}
If \(N=p^a\), the PD-principality calculation above gives
\[
 M\cong U_R(\mathfrak a)
 \quad\hbox{freely on }T^{[N]},\qquad
 M_{\mathrm{sp}}\cong A.
 \tag{CON.10}
\]
To verify this with the operator action, put
\(C=\Gamma_R(U)\), \(U\mapsto T^{[N]}\), and \(J=C_+\).
Since \(S\) is \(C\)-free and \(I=SJ\), one has
\(I/I^2=S\otimes_C(J/J^2)\).
The abelianized free PD algebra \(J/J^2\) is the free
\(B\)-module on \(U\) (Kmail--Kozak--Miller, Corollary~6.5).
Thus \(I/I^2=A\otimes_R B=U_R(\mathfrak a)\).
The tensor product PD formula gives exactly the operator action
in (CON.1), so this is an isomorphism of Beck modules. Quotienting
the operator directions gives (CON.10).

\paragraph{Consequence for the current gate.}
The coefficient tail has genuinely different linearizations.
Ordinary \(S\)-module projective dimension is infinite even for a
PD-principal tail. At a prime \(p>2N\), Beck conormal projective
dimension is also infinite, but the special conormal relevant to
\(d\gamma_r(x)=\gamma_{r-1}(x)\,dx\) has projective dimension one.
These facts neither compute higher derived special cotangent groups
nor provide coherent multiplicative and simplicial choices of the
repair primitives. They exclude an inference from Beck infinite
projective dimension to failure of special dg-PD completion.

\paragraph{Source locator and check.}
The operator description used here is from A.~Kmail, J.~Kozak and
H.~Miller, \emph{Divided powers and K\"ahler differentials},
Proposition~4.2, Theorem~5.4 and Corollary~6.5; author PDF
\url{https://math.mit.edu/~hrm/papers/dp5.pdf}.
The tail calculation and all changes of scalars above were read
in full by the author and independently by cote~15, GPT-6 Astra Pro,
effort not exposed, Europe/Berlin date 2026-11-22 (PASS).
Root's check is recorded separately.

